\section{Referencia rápida de terminología y patrones de diseño}

\subsection{Comparación de términos clave}
Para evitar desviaciones semánticas cuando los equipos colaboran entre módulos, esta clase permite extraer un glosario mínimo. Se recomienda fijar estas definiciones de términos en la documentación del proyecto, para reducir el costo de comunicación y las desviaciones experimentales causadas por «mismo nombre, distinto significado».

\begin{table}[H]
\centering
\begin{tabular}{p{2.8cm}p{3.6cm}p{4.1cm}}
\toprule
\textbf{Término} & \textbf{Significado en el contexto del curso} & \textbf{Puntos a cuidar al implementarlo} \\
\midrule
Reasoning action & Paso de pensamiento intermedio usado para mejorar la calidad de las acciones posteriores & Debe vincularse a una condición de activación y a un límite de presupuesto \\
Reflection & Autorrevisión de la trayectoria existente con intento de corrección & Requiere una señal de fallo; de lo contrario tiende a girar en vacío \\
Replanning & Reescritura del plan debido a una desviación del entorno & Debe definirse con claridad el umbral de replanificación, para evitar oscilaciones frecuentes \\
Working memory & Contexto de corto plazo necesario para el paso actual & Controlar su capacidad, para evitar que desplace información clave \\
Episodic memory & Registro de eventos y resultados durante el proceso de la tarea & Debe incluir origen y marca de tiempo, para permitir la revisión posterior \\
Semantic memory & Contenido estable de hechos y base de conocimiento & Gestionar versiones y conflictos, para evitar la contaminación con conocimiento obsoleto \\
Procedural memory & Estrategias, procesos y plantillas reutilizables & Mantener las condiciones de activación, para evitar generalizaciones erróneas \\
Verifier & Mecanismo que verifica la corrección intermedia o final & El propio verificador también requiere evaluación continua \\
\bottomrule
\end{tabular}
\caption{Explicación de ingeniería de los términos relacionados con Reasoning, Memory y Planning}
\end{table}

\subsection{Patrones de diseño frecuentes}
La práctica del curso puede condensarse además en varios patrones de diseño de alta frecuencia. No son soluciones universales, pero reducen de manera considerable la fricción de migrar de un prototipo de investigación a un sistema de producción.

\begin{knowledgebox}{Cuatro patrones de diseño de uso común}
\begin{enumerate}
    \item \textbf{Plan-Execute-Verify}: primero planificar, después ejecutar y luego verificar; adecuado para tareas que se pueden verificar;
    \item \textbf{Think-Act-Reflect}: alternancia entre razonamiento y ejecución; adecuado para entornos dinámicos;
    \item \textbf{Retrieve-Reason-Write}: primero obtener evidencia y después razonar; adecuado para tareas intensivas en conocimiento;
    \item \textbf{Shadow-to-Production}: primero evaluación en modo sombra y después despliegue gradual; adecuado para negocios de alto riesgo.
\end{enumerate}
\end{knowledgebox}

\begin{warningbox}{Condiciones límite al reutilizar un patrón}
El efecto del mismo patrón puede ser opuesto en distintas distribuciones de tareas. Antes de ponerlo en producción debe verificarse si se cumplen las premisas del patrón: que la tarea sea resoluble, la estabilidad de las herramientas, el presupuesto de latencia y las restricciones de permisos.
\end{warningbox}

\subsection{Resumen de la sección}
Unificar la terminología y los patrones de diseño reduce la fricción de colaboración entre equipos y aumenta la comparabilidad de los experimentos y la mantenibilidad del sistema. Este es también un paso clave para que la ingeniería de Agent pase de la ``experiencia individual'' a la ``capacidad organizacional''.

